% Sections 3.1 to 3.6, from lectures/L03/appendix/a_theory.md.

\section{Ekvationer}
\label{c3:sec:ekvationer}
En \textbf{ekvation} är ett påstående om att två uttryck är lika:
\[
  \text{vänsterled} = \text{högerled}
\]
En \textbf{linjär ekvation} i en obekant $x$ är en ekvation där $x$ förekommer med högst potensen
1:
\[
  ax + b = c, \quad a \neq 0
\]
Att \textbf{lösa en ekvation} innebär att bestämma vilket eller vilka värden på $x$ som uppfyller
ekvationen.

\section{Ekvationsprinciperna}
\label{c3:sec:principer}
En ekvations lösning förändras inte om man utför \emph{samma} operation på båda led.

\begin{booktable}{@{}p{12.5em}L@{}}
  \thead{Princip} & \thead{Beskrivning} \\
  \midrule
  \textbf{Additionsprincipen} & Addera eller subtrahera samma tal i båda led \\
  \textbf{Multiplikationsprincipen} & Multiplicera eller dividera med samma tal ($\neq 0$) i båda
    led \\
\end{booktable}

\textbf{Strategi:} Flytta alla termer med obekant till ett led och konstanterna till det andra.

\section{Lösa linjära ekvationer steg för steg}
\label{c3:sec:losa}

\begin{exempel}
Lös $3x - 5 = 7$.
\begin{alignat*}{2}
  3x - 5 &= 7 \\
  3x &= 12 &\qquad& \text{addera $5$ i båda led} \\
  x &= 4 && \text{dividera med $3$ i båda led}
\end{alignat*}
Kontroll: $3 \cdot 4 - 5 = 12 - 5 = 7$, vilket stämmer.
\end{exempel}

\begin{exempel}[Ohms lag]
En krets har spänningen $U = 12\,\text{V}$ och strömmen $I = 0{,}5\,\text{A}$. Beräkna motståndet
$R$. Ohms lag ger
\[
  U = R \cdot I \quad \Leftrightarrow \quad 12 = R \cdot 0{,}5.
\]
Dividera med $0{,}5$:
\[
  R = \frac{12}{0{,}5} = 24\,\Omega
\]
\end{exempel}

\begin{exempel}[Med parenteser]
Lös $2(3x + 1) = 14$.
\begin{alignat*}{2}
  6x + 2 &= 14 &\qquad& \text{multiplicera ut parentesen} \\
  6x &= 12 && \text{subtrahera $2$} \\
  x &= 2 && \text{dividera med $6$}
\end{alignat*}
\end{exempel}

\section{Proportionsekvationer}
\label{c3:sec:proportion}
En proportionsekvation har formen
\[
  \frac{a}{b} = \frac{c}{d}
\]
och löses genom \textbf{korsvis multiplikation}, det vill säga genom att multiplicera diagonalt:
\[
  ad = bc
\]

\begin{exempel}
Lös $\dfrac{x}{6} = \dfrac{4}{3}$.
\[
  3x = 24 \quad \Rightarrow \quad x = 8
\]
\end{exempel}

\section{Olikheter}
\label{c3:sec:olikheter}
En \textbf{olikhet} anger att ett uttryck är större eller mindre än ett annat:

\begin{narrowtable}{@{}ll@{}}
  \thead{Symbol} & \thead{Läses som} \\
  \midrule
  $a < b$ & $a$ är mindre än $b$ \\
  $a > b$ & $a$ är större än $b$ \\
  $a \leq b$ & $a$ är mindre än eller lika med $b$ \\
  $a \geq b$ & $a$ är större än eller lika med $b$ \\
\end{narrowtable}

\subsection{Lösa olikheter}
\label{c3:sec:losaolikheter}
Olikheter löses precis som ekvationer, med ett viktigt undantag:

\begin{aside}
\note[Obs] Om man multiplicerar eller dividerar med ett \textbf{negativt tal} vänder
olikhetstecknet.
\end{aside}

\begin{exempel}
Lös $2x + 3 > 9$.
\[
  2x > 6 \quad \Rightarrow \quad x > 3
\]
Lösningen är alla $x$ som är större än $3$, det vill säga $x \in (3, \infty)$.
\end{exempel}

\begin{exempel}[Negativt tal]
Lös $-3x \leq 12$. Dividera med $-3$; tecknet vänder:
\[
  x \geq -4
\]
\end{exempel}

\subsection{Praktisk tillämpning: säkert driftområde}
\label{c3:sec:driftomrade}
En komponent tål högst $P_{\max} = 2\,\text{W}$, och effekten ges av $P = U \cdot I$ med
$I = 0{,}5\,\text{A}$. Det tillåtna spänningsintervallet bestäms av
\[
  U \cdot 0{,}5 \leq 2 \quad \Rightarrow \quad U \leq 4\,\text{V}.
\]
Spänningen får vara högst $4\,\text{V}$ för att komponenten inte ska skadas.

\section{Absolutbelopp i ekvationer}
\label{c3:sec:absolutbelopp}
Ekvationer med absolutbelopp har i allmänhet två lösningar. Ekvationen $\abs{ax + b} = c$, med
$c \geq 0$, ger
\[
  ax + b = c \quad \text{eller} \quad ax + b = -c.
\]

\begin{exempel}
Lös $\abs{2x - 3} = 5$.
\begin{align*}
  \text{Fall 1:} &\quad 2x - 3 = 5 \quad \Rightarrow \quad x = 4 \\
  \text{Fall 2:} &\quad 2x - 3 = -5 \quad \Rightarrow \quad x = -1
\end{align*}
Lösningen är $x = 4$ eller $x = -1$.
\end{exempel}
